\section{Independent hierarchical messages and the cavity comparison}
\label{comp:section}

We turn the reservoir into a trial law in three steps: encode its full
one-site branching law, use multioverlap concentration to separate the
sites, and compare the resulting insertion logarithms.

Fix $\eps>0$. At each depth $L\ge2$, use the sequence of sizes and
deterministic parameter choices supplied by \cref{pert:selection}.
Write $N$ for sizes on this sequence. The choices may depend on $L$, and all statements below concern
these choices. Retain $\underline F$ for the liminf over \emph{all} system
sizes of the bounded model, as before selection. In particular,
\begin{equation}\label{comp:low-increments}
 \Delta_N^{L,\mathbf u_{L,N}}\le\underline F+2\eps.
\end{equation}

\subsection{Encoding a single-site tree}

\begin{lemma}[Recursive law encoding]\label[lemma]{comp:encoding}
For each fixed $N$, $L\ge2$, and selected parameter vector, with the reservoir
root disorder averaged, there is a deterministic trial law
$\zeta_{L,N}\in\mathcal M_L$ whose single-label spin process has the same
expected spin products on every prescribed tree as a uniformly chosen site
of the reservoir. Different trial labels can be chosen independently.
\end{lemma}

\begin{proof}
Include an independent uniform site $i\in[N]$ in the root variable, along
with the reservoir root disorder. Sample the tilted reservoir auxiliary
kernels along each path, and let $\mathscr G_d$ be the sigma-field generated
by this enlarged root and the full auxiliary prefix through depth $d$.
At depth $r=L-1$ put
\[
 x=\arctanh\angles{\sigma_i},
\]
where the brackets denote the leaf Gibbs expectation. This is finite:
at finite $N$ all configuration weights are strictly positive, so
$-1<\angles{\sigma_i}<1$.
Thus $x$ is determined at depth $r$, before the terminal spin transition
at depth $L$. At each terminal replica leaf, draw a spin with law $q_x$. Given the
auxiliary prefix, these spins are independent and have exactly the
one-site law of the independent terminal Gibbs replicas.

Define the random measures backward by
\[
 \eta_{r-1}=\operatorname{Law}(x\mid\mathscr G_{r-1}),\qquad
 \eta_d=\operatorname{Law}(\eta_{d+1}\mid\mathscr G_d)
 \quad(0\le d<r-1).
\]
Then $\eta_d\in\mathcal M_{r-d}$, so
$\zeta_{L,N}=\operatorname{Law}(\eta_0)\in\mathcal M_L$.
All state spaces are standard
Borel, so these conditional probability kernels have measurable versions.
The auxiliary priors used here may be taken to be finite products of finite
mark spaces at each realized root; hence the same assertion also follows
directly by finite sums conditional on the root counts.

Here is an explicit verification for every branching pattern. Fix a
descendant tree $T$ and prescribe a spin $e_a$ at each terminal leaf.
At depth $r$, its conditional probability is
$G_T(x)=\prod_a q_x(e_a)$. At depth $d<r$, let $T_1,\ldots,T_j$ be
the subtrees rooted at its children. Conditional independence of children
gives the backward recursion
\[
 G_T(\eta_d)=\prod_{\ell=1}^j
       \int G_{T_\ell}(\nu)\,\eta_d(d\nu).
\]
By induction, this probability depends on the full prefix only through
$\eta_d$. Thus the nested measures preserve the joint output law of every
descendant tree, including unary stretches and arbitrary numbers of
children. This recursion is precisely the branching version of
\eqref{model:chain}. Taking $\zeta_{L,N}=\operatorname{Law}(\eta_0)$
preserves the root mixture as well. In particular, for every subset $S$
of the leaves,
\begin{equation}\label{comp:single-label-moments}
 \E\prod_{a\in S}\tau^a
 =\frac1N\sum_{i=1}^N\E\prod_{a\in S}\sigma_i^a
 =\E R_S.
\end{equation}
Use independent copies of the entire construction, including the reservoir
root and uniform site, for different labels. The resulting messages are
independent of the new physical disorder in the cavity insertions.
\end{proof}

\subsection{From multioverlaps to spin patterns}

\begin{lemma}[Empirical-pattern comparison]\label[lemma]{comp:patterns}
Fix a prescribed tree with $k$ leaves. Suppose an observable of absolute
value at most one uses the replica-spin patterns at $a$ independent uniform
site labels. Replacing these patterns by $a$ independent trial-label patterns
from \cref{comp:encoding} changes its expected value by at most
\begin{equation}\label{comp:pattern-error}
 a\sum_{\varnothing\ne S\subseteq[k]}\E|R_S-\E R_S|.
\end{equation}
The observable may also depend on new insertion disorder, independent of
the reservoir.
\end{lemma}

\begin{proof}
Conditional on the sampled reservoir configurations, the pattern
$\boldsymbol\sigma_i=(\sigma_i^1,\ldots,\sigma_i^k)$ at a uniform site
has empirical distribution
\[
 \pi(\boldsymbol e)=\frac1N\sum_i
 \prod_{b=1}^k\frac{1+e_b\sigma_i^b}{2}
 =2^{-k}\sum_{S\subseteq[k]}R_S\prod_{b\in S}e_b,
 \qquad \boldsymbol e\in\{-1,1\}^k.
\]
The trial-label distribution is $\bar\pi=\E\pi$ by
\eqref{comp:single-label-moments}. Independent uniform sampling with
replacement gives $\pi^{\otimes a}$ conditionally in the reservoir; the
trial distribution is $\bar\pi^{\otimes a}$. The elementary telescoping
bound for product measures gives
\[
 \norm{\pi^{\otimes a}-\bar\pi^{\otimes a}}_{\rm TV}
 \le a\norm{\pi-\bar\pi}_{\rm TV}
 \le\frac a2\sum_{\varnothing\ne S\subseteq[k]}|R_S-\E R_S|.
\]
Integrating a function bounded by one costs at most twice total variation.
Conditioning on the independent insertion disorder proves the last assertion.
\end{proof}

Let $Q_k$ be the tree law in \cref{tree:log-expansion}. We need the following
consequence of the decorrelation theorem:
\begin{equation}\label{comp:shape-error}
 \lim_{L\to\infty}\limsup_{N\to\infty}
 \E_{Q_k}\sum_{\varnothing\ne S\subseteq[k]}
 \E|R_S-\E R_S|=0\qquad(k\text{ fixed}).
\end{equation}
Here and below the inner limit is along the selected fixed-$L$ sequence.
To see this, fix $0<\eta<1/2$, and let $\mathfrak S_k$ be the finite collection
of labeled topologies on nonempty subsets of $[k]$. We claim the bound
\begin{equation}\label{comp:shape-error-bound}
 \E_{Q_k}\sum_{\varnothing\ne S\subseteq[k]}
       \E|R_S-\E R_S|
 \le C_k\left(\eta+L^{-1}
       +\sum_{U\in\mathfrak S_k}
          \bigl(\mathcal A_{L,\eta}\Var(R_U)\bigr)^{1/2}\right).
\end{equation}
The average for a singleton has no depth coordinates.
By \cref{tree:regular-shapes}, nonregular trees
have probability at most $C_k(\eta+L^{-1})$. On regular trees every induced
subtree retains the interior and separation conditions on its branching
depths. For a topology with $a$ branching vertices, each depth assignment
has probability at most $C_kL^{-a}$. If an induced subtree retains $a'$
branching vertices, each assignment of its depths has at most $L^{a-a'}$
extensions to the full tree. Summing out these unused coordinates leaves
exactly the normalization $L^{-a'}$ in $\mathcal A_{L,\eta}$.
Cauchy--Schwarz bounds the resulting average of mean absolute deviations
by $\bigl(\mathcal A_{L,\eta}\Var(R_U)\bigr)^{1/2}$.
Summing over the finitely many topologies and subsets, and using the bound
two for each deviation on nonregular trees, proves
\eqref{comp:shape-error-bound}. Apply \cref{decor:overlap-concentration}
first at fixed $\eta$, taking the fixed-$L$ limit in $N$ before
$L\to\infty$, and then let $\eta\downarrow0$ to obtain
\eqref{comp:shape-error}.

\subsection{Comparison of the insertion logarithms}

\begin{proposition}\label[proposition]{comp:insertion-comparison}
For the trial laws of \cref{comp:encoding}, with
$\mathbf m^{(L)}=(1/L,\ldots,(L-1)/L)$,
\begin{equation}\label{comp:increment-comparison}
 \lim_{L\to\infty}\limsup_{N\to\infty}
 \left|\Delta_N^{L,\mathbf u_{L,N}}
 -\cB_{L-1}(\zeta_{L,N},\mathbf m^{(L)})\right|=0.
\end{equation}
\end{proposition}

\begin{proof}
For each insertion in \cref{cavity:increment}, replace the reservoir spins
at its site labels by independent spin processes from \cref{comp:encoding}.
Denote the resulting trial insertions by $\widetilde A_{\rm s}$ and
$\widetilde A_{\rm b}$, and their recursions from
\eqref{tree:Y-recursion} by $\widetilde Y_{0,\rm s}$ and
$\widetilde Y_{0,\rm b}$. These recursions use the trial kernels.

At the terminal spin transition, averaging $\widetilde A_{\rm s}$ over
its independent trial spins gives exactly $V_{\rm s}$ in
\eqref{model:vs}; hence $\widetilde Y_{0,\rm s}=T_0V_{\rm s}$.
For $\widetilde A_{\rm b}$ the same transition gives a product of
$K_{\rm b}$ edge factors with disjoint message labels. Given the physical
disorder and all message prefixes at a depth, these label groups remain
independent. Their product therefore factors through each conditional
power mean, successively from the last level to the first. Consequently,
\[
 \begin{aligned}
 \E\log\widetilde Y_{0,\rm b}
 &=\alpha(p-1)\E\log T_0V_{\rm e},\\
 \log2+\E\log\widetilde Y_{0,\rm s}
       -\E\log\widetilde Y_{0,\rm b}
 &=\cB_{L-1}(\zeta_{L,N},\mathbf m^{(L)}).
 \end{aligned}
\]
By \eqref{cavity:identity}, it remains to compare the reservoir and trial
insertion logarithms in the stated limits.

Consider either insertion and first restrict its
Poisson count $K$ to $K\le M$, where $M$ is fixed. Set
$C=H+BM$, $z=\cosh C$, and $c=\tanh C<1$. On this event both the reservoir
insertion and its trial version can be written as
$A=z(1+D)$ with $|D|\le c$, uniformly in all disorder and path variables.
The deterministic term $\log z$ cancels in the comparison.

At each order $k$ of \cref{tree:log-expansion}, the product of the $k$
copies of $D$ is a bounded function of replica-spin patterns at at most
$pM$ labels and of independent insertion disorder. By
\cref{comp:patterns}, its evaluation error on a fixed tree is bounded by
$pM\sum_{S\ne\varnothing}\E|R_S-\E R_S|$. Hence
\eqref{comp:shape-error} makes its tree-averaged error tend to zero in the
ordered limits. For every fixed $M$, truncating the expansion after $q$
terms leaves an error at most
$\sum_{k>q}c^k/k$ on either side, uniformly in $N,L$.

Outside the count cutoff the original insertion logarithms, on both sides,
are bounded by $H+BK$. Their expected contribution is therefore at most
$2\E[(H+BK)\mathbf1_{\{K>M\}}]$, which tends to zero as $M\to\infty$.
This proves convergence of the expected insertion logarithms: explicitly,
first choose $M$ to control the Poisson tail, then $q$ to control the
geometric tail, and then take the limits in $N$ and $L$ for the finitely
many remaining expansion terms. Together with the trial identities above
and \eqref{cavity:identity}, this proves \eqref{comp:increment-comparison}.
\end{proof}

\begin{corollary}\label[corollary]{comp:bounded-lower}
For every bounded permutation-invariant interaction law and bounded field
law, with the functionals defined in \cref{model:section},
\[
 \liminf_{N\to\infty}F_N\ge\inf_{r\ge0}\Phi_r.
\]
\end{corollary}

\begin{proof}
Write $A=\inf_{r\ge0}\Phi_r$, which is finite by
\eqref{model:trial-bounds}. Each encoded trial value is at least $A$, so
\eqref{comp:low-increments} gives, at every fixed $L$,
\[
 A\le\underline F+2\eps+
 \limsup_{N\to\infty}
 \left|\Delta_N^{L,\mathbf u_{L,N}}
       -\cB_{L-1}(\zeta_{L,N},\mathbf m^{(L)})\right|.
\]
The inner limit uses that depth's selected sequence. Letting $L\to\infty$
in \cref{comp:insertion-comparison} gives
$A\le\underline F+2\eps$, and then $\eps\downarrow0$ proves the result.
\end{proof}
